% Options for packages loaded elsewhere
\PassOptionsToPackage{unicode}{hyperref}
\PassOptionsToPackage{hyphens}{url}
%
\documentclass[
  english,
]{article}
\usepackage{lmodern}
\usepackage{amssymb,amsmath}
\usepackage{ifxetex,ifluatex}
\ifnum 0\ifxetex 1\fi\ifluatex 1\fi=0 % if pdftex
  \usepackage[T1]{fontenc}
  \usepackage[utf8]{inputenc}
  \usepackage{textcomp} % provide euro and other symbols
\else % if luatex or xetex
  \usepackage{unicode-math}
  \defaultfontfeatures{Scale=MatchLowercase}
  \defaultfontfeatures[\rmfamily]{Ligatures=TeX,Scale=1}
\fi
% Use upquote if available, for straight quotes in verbatim environments
\IfFileExists{upquote.sty}{\usepackage{upquote}}{}
\IfFileExists{microtype.sty}{% use microtype if available
  \usepackage[]{microtype}
  \UseMicrotypeSet[protrusion]{basicmath} % disable protrusion for tt fonts
}{}
\makeatletter
\@ifundefined{KOMAClassName}{% if non-KOMA class
  \IfFileExists{parskip.sty}{%
    \usepackage{parskip}
  }{% else
    \setlength{\parindent}{0pt}
    \setlength{\parskip}{6pt plus 2pt minus 1pt}}
}{% if KOMA class
  \KOMAoptions{parskip=half}}
\makeatother
\usepackage{xcolor}
\IfFileExists{xurl.sty}{\usepackage{xurl}}{} % add URL line breaks if available
\IfFileExists{bookmark.sty}{\usepackage{bookmark}}{\usepackage{hyperref}}
\hypersetup{
  pdflang={en},
  hidelinks,
  pdfcreator={LaTeX via pandoc}}
\urlstyle{same} % disable monospaced font for URLs
\usepackage[margin=1in]{geometry}
\usepackage{longtable,booktabs}
% Correct order of tables after \paragraph or \subparagraph
\usepackage{etoolbox}
\makeatletter
\patchcmd\longtable{\par}{\if@noskipsec\mbox{}\fi\par}{}{}
\makeatother
% Allow footnotes in longtable head/foot
\IfFileExists{footnotehyper.sty}{\usepackage{footnotehyper}}{\usepackage{footnote}}
\makesavenoteenv{longtable}
\setlength{\emergencystretch}{3em} % prevent overfull lines
\providecommand{\tightlist}{%
  \setlength{\itemsep}{0pt}\setlength{\parskip}{0pt}}
\setcounter{secnumdepth}{-\maxdimen} % remove section numbering
\ifxetex
  % Load polyglossia as late as possible: uses bidi with RTL langages (e.g. Hebrew, Arabic)
  \usepackage{polyglossia}
  \setmainlanguage[]{english}
\else
  \usepackage[shorthands=off,main=english]{babel}
\fi

\author{}
\date{}

\begin{document}

\hypertarget{pointwise-einstein-approximation-exact-models-and-closed-gluing}{%
\section{Pointwise Einstein approximation, exact models, and closed
gluing}\label{pointwise-einstein-approximation-exact-models-and-closed-gluing}}

\hypertarget{the-construction-problem-and-the-inherited-results}{%
\subsection{1. The construction problem and the inherited
results}\label{the-construction-problem-and-the-inherited-results}}

The relevant error is the unweighted endomorphism norm \[
\delta_\sigma(g)=\|g^{-1}\operatorname{Ric}_g-\sigma I\|_{L^\infty,\mathrm{op}},
\qquad \sigma\in\{-1,1\}.
\] On a fixed closed smooth manifold, constructing
\(\delta_\sigma(g_j)\to0\) establishes \(e_\sigma(M)=0\). It does not
decide whether another metric on \(M\) is Einstein. An exact metric on
an open model, a family on a fixed closed manifold, and an obstruction
to repairs in a specified neighborhood are different results. This
chapter develops their connections.

The baseline already contains two complete approximation mechanisms. The
first resolves a fixed Einstein four-orbifold with marked \(A_1\)
fillings. A corrected Eguchi--Hanson core has the prescribed nonzero
Einstein source. In a common radial gauge, its quadratic mismatch with
the exterior is a pure Weyl tensor \[
K_{ij}=-\tfrac13 C_{ikjl}x^kx^l,\qquad
\operatorname{tr}K=\operatorname{div}K=\Delta K=0.
\] For \(s=\log r\), \[
D\operatorname{Ric}_e[\eta(s)K]
 =-\frac{\eta''+6\eta'}{2r^2}K.
\] Since \(K=O(r^2)\), interpolation over a logarithmic interval of
length \(L\) costs \(O(L^{-1})\) in the required norm. The core,
annulus, and exterior estimates give
\(\delta_\sigma(g_a)\le C/|\log a|\) on the fixed marked resolution.
Core correction and nonlinear tail hypotheses are part of that theorem,
not consequences of the linear identity. See the
\href{../baseline/research/20260916_construction.tex}{baseline
construction manuscript, §§2--5}.

The second mechanism fixes a surface bundle \(E_D\) with Dehn-twist
monodromy and interpolates hyperbolic fibre metrics along a long
hyperbolic base collar. The complete Ricci matrix has error
\(O(L^{-1})\); the shrinking angular direction creates no loss because
coefficients are angularly invariant. The baseline also permits
monodromy through a geometric free-group quotient, including
noncommuting mapping classes. Its balanced double-collar family has
fibre waist \(\varepsilon\), base waist \(\beta=\varepsilon^2\), and \[
\delta_-(G_\varepsilon)=\frac{2}{\pi^4}\varepsilon^2+O(\varepsilon^4).
\] An \(O(\varepsilon^2)\) global correction reduces the defect to
\(O(\varepsilon^3)\). These are inherited constructions, not results
first obtained in the subsequent twelve rounds. See
\href{../baseline/research/20260912.tex}{the detailed bundle chapter}
and \href{../baseline/research/20260916_global.tex}{the baseline global
manuscript, section bundle-input}.

The subsequent work has four centers: a complete exact double-cylinder
family; nonlinear small-repair obstructions on closed \(E_D\); complete
two-ended approximate metrics with fixed circle topology; and closed
sheared product-block metrics beyond surface bundles. No
publication-priority claim is made for the constructions or methods.

\hypertarget{the-complete-invariant-double-cylinder-problem}{%
\subsection{2. The complete invariant double-cylinder
problem}\label{the-complete-invariant-double-cylinder-problem}}

\hypertarget{all-bounded-invariant-jacobi-fields}{%
\subsubsection{2.1. All bounded invariant Jacobi
fields}\label{all-bounded-invariant-jacobi-fields}}

On \(X=(\mathbb R\times S^1)^2\), fix angular periods and put \[
g_0=dt^2+\cosh^2t\,dy^2+du^2+\cosh^2u\,dx^2,\quad
f=\operatorname{sech}^2t,\quad r=\operatorname{sech}^2u,\quad
J(v)=\arctan(\sinh v),\quad W=\cosh t\cosh u.
\] Use the orthonormal coframe
\((e^0,e^1,e^2,e^3)=(dt,\cosh t\,dy,du,\cosh u\,dx)\). Let
\(T_1=(e^0)^2-(e^1)^2\), \(S_1=e^0e^1+e^1e^0\), and similarly
\(T_2,S_2\) on the second factor. The Bianchi-gauged operator is
\(L=(\nabla^*\nabla-2\mathring R)/2\). Writing
\(P_t=\partial_t^2+\tanh t\,\partial_t\), its ten tensor blocks are

\begin{longtable}[]{@{}ll@{}}
\toprule
Block & Operator \(2L\)\tabularnewline
\midrule
\endhead
Two factor traces & \(-P_t-P_u+2\)\tabularnewline
\(T_1,S_1\) & \(-P_t-P_u+2-4f\)\tabularnewline
\(T_2,S_2\) & \(-P_t-P_u+2-4r\)\tabularnewline
Four mixed polarizations & \(-P_t-P_u+2-f-r\)\tabularnewline
\bottomrule
\end{longtable}

Every bounded angularly invariant Jacobi field, initially allowing
arbitrary joint \((t,u)\)-dependence, is \[
H[z]=f[(a_0+a_1J(u))T_1+(b_0+b_1J(u))S_1]
 +r[(A_0+A_1J(t))T_2+(B_0+B_1J(t))S_2].
\tag{2.1}
\] The classification is not a separated-variables assumption. The block
estimates exclude the remaining bounded modes. On the full invariant
\(L^2(W\,dt\,du)\) domain there is instead a positive gap: \[
\langle Lv,v\rangle\ge\tfrac18\|v\|_2^2.
\tag{2.2}
\] The eight bounded fields in (2.1) are not \(L^2\) on the whole
product. See \href{../proofs/round2_construction.md}{the kernel and gap
proof} and its
\href{../proofs/round2_construction_review.md}{independent review}.

\hypertarget{integrating-the-entire-nonlinear-source}{%
\subsubsection{2.2. Integrating the entire nonlinear
source}\label{integrating-the-entire-nonlinear-source}}

For constant parameters define a complete curvature-\(-1\) cylinder \[
h_{\alpha,\gamma}
=\frac{\cosh^2v}{\sinh^2v+e^{-\alpha}}\,dv^2
+(\sinh^2v+e^{-\alpha})(d\theta+\gamma\operatorname{sech}^3v\,dv)^2.
\] Let the first parameters depend on \(u\) through \(a_0+a_1J(u)\),
\(b_0+b_1J(u)\), and the second pair analogously on \(t\). Their sum
defines \(g[z]\), with first variation \(H[z]\) and determinant one in
the background frame.

\textbf{Exact model theorem.} For sufficiently small
\(z\in\mathbb R^8\), there is an analytic correction \(k[z]\) satisfying
\[
G[z]=g[z]+k[z],\qquad \operatorname{Ric}_{G[z]}=-G[z],\qquad
\beta_{g[z]}k[z]=0,\qquad
\|k[z]\|_{H^{m+2}_{1/4}}\le C|z|^2
\tag{2.3}
\] for \(m=4\), with higher regularity and
\(|\nabla^jk[z]|\le C_j|z|^2W^{-3/4}\) for every fixed \(j\). The
metrics are complete, smooth, and invariant under both rotations.
Uniqueness is in the stated gauge and correction space, not modulo
arbitrary diffeomorphisms. Eight distinct moduli are not asserted.

Here \(H^m_\delta\) means \(W^\delta\) times the tensor belongs to the
background invariant Sobolev space. The analytic proof has five steps.

\begin{enumerate}
\def\labelenumi{\arabic{enumi}.}
\tightlist
\item
  Compactify by \(\xi=J(t)\), \(\eta=J(u)\). All frame coefficients
  extend smoothly to the closed square, with radial derivatives
  \(\cos\xi\,\partial_\xi\), \(\cos\eta\,\partial_\eta\).
\item
  The exact nonlinear residual vanishes on every face and corner.
  Hadamard division gives \(E[z]=(\cos\xi\cos\eta)R[z]\), with
  \(R[z]=O(|z|^2)\) smooth. Every nonlinear order therefore lies in
  \(H^m_\delta\), \(\delta<1/2\).
\item
  Conjugate \(L\) by \(W^\delta\). Its drift and linear zeroth-order
  terms are skew-adjoint together; \(|\nabla\log W|^2\le2\) gives
  \(\operatorname{Re}\langle L_\delta v,v\rangle \ge(1/8-\delta^2)\|v\|^2\).
  At \(\delta=1/4\), Lax--Milgram and regularity give the full inverse.
\item
  Solve
  \(\operatorname{Ric}(g[z]+k)+g[z]+k+ \delta^*_{g[z]+k}\beta_{g[z]}k=0\)
  by contraction in a ball of radius \(C|z|^2\).
\item
  The \(L^2\) Bianchi one-form obeys
  \((\nabla^*\nabla-\operatorname{Ric})\omega=0\). Negative Ricci
  already follows from \(C^2\) proximity; cutoff integration forces
  \(\omega=0\).
\end{enumerate}

A moving reference is essential. A constant cylinder length change has a
fixed-background Bianchi defect independent of the other radial
coordinate, hence not \(L^2\); its quadratic size does not fix that. The
\href{../proofs/round3_construction.md}{nonlinear proof} and
\href{../proofs/round3_construction_review.md}{review} check the whole
source, spaces, inverse, positivity, and gauge removal. The general
deformation method is established prior work; the specified degenerate
cylinder factors and end data are treated directly here.

\hypertarget{nonlinear-charge-boundary-flux-and-exact-insertion}{%
\subsection{3. Nonlinear charge, boundary flux, and exact
insertion}\label{nonlinear-charge-boundary-flux-and-exact-insertion}}

Write the exact metric in mechanical-connection form \[
G=h+H_{ab}(d\theta^a+\mathcal A^a)(d\theta^b+\mathcal A^b),\quad
v=\sqrt{\det H},\quad F=d\mathcal A,\quad q=vH*_hF.
\] Use angular generators \((\partial_y,\partial_x)\) and quotient
orientation \(dt\wedge du\). The mixed Einstein equation gives \(dq=0\).
Neither \(H\) nor \(h\) is assumed diagonal. Set \(c=\pi^2/2\) and \[
M[z]=\int_{\mathbb R^2}v^{-1}H^{-1}\,dA_h,\qquad
S[z]=c(-b_1,B_1)^T.
\] Then \[
\int_{\mathbb R^2}F=S[z]=M[z]q[z],\qquad M[z]>0,\qquad M[0]=cI.
\tag{3.1}
\] The matrix density is dominated by
\(\operatorname{diag}(\operatorname{sech}^3t\operatorname{sech}u, \operatorname{sech}t\operatorname{sech}^3u)\).
The correction changes the connection by a form whose boundary integral
over expanding squares tends to zero, also after fixed parameter
derivatives. Stokes computes the total curvature from actual end
markings.

The finite nonlinear twist energy is \[
\mathcal E_{\rm tw}=q^TMq=q\cdot S=S^TM^{-1}S
=\lim_{R\to\infty}\int_{\partial[-R,R]^2}q\cdot\mathcal A.
\tag{3.2}
\] It vanishes exactly when \(b_1=B_1=0\). Constant twists \(b_0,B_0\)
are removed by angular shifts using
\(\int_0^v\operatorname{sech}^3w\,dw\); they are exact gauge in this
family, though end markings change. These identities are not a full
metric Dirichlet-to-Neumann operator. See
\href{../proofs/round4_construction.md}{the charge and energy proof}.

The quadratic two-face calculation remains useful. Its net sources are
\[
I_1=\tfrac{\pi^2}{8}(a_1^2-3b_1^2-A_1^2-B_1^2),\quad
I_2=\tfrac{\pi^2}{8}(A_1^2-3B_1^2-a_1^2-b_1^2).
\] Homogeneous second-order end shifts change opposite face fluxes but
not their sum \(-I_i\). Vanishing net exterior response would imply
\(b_1=B_1=0\), \(a_1^2=A_1^2\). That hypothesis is not known on \(E_D\).
The true exterior equation retains its source and the detector's adjoint
defect; physical periods multiply every source and flux by the same
factor. Local convergence alone does not discard these normalized terms.
See \href{../proofs/round3_matching_next.md}{the conditional matching
law}.

There is an unconditional exact-insertion obstruction. If
\((b_1,B_1)\ne0\), no nonempty open piece of \(G[z]\) is locally
isometric to an open piece of a connected closed Einstein four-manifold.
Its complete universal cover has fixed-radius balls escaping to product
geometry, with \(\nabla\mathrm{Rm}\to0\). A cocompact isometry group
would translate them into a compact set, force \(\nabla\mathrm{Rm}=0\)
on an open ball, then everywhere by analyticity. The cover would be
\(\mathbb H^2\times\mathbb H^2\). Every locally free two-dimensional
abelian Killing action there has integrable orthogonal complement, so
\(F=0\), contradicting (3.1). Analytic continuation of a local isometry
between complete simply connected analytic covers transfers any proposed
closed insertion to this contradiction. The
\href{../proofs/round4_exact_core.md}{proof} excludes unchanged exact
cores, not closed metrics converging to the model or corrections on
every open set.

\hypertarget{a-nonlinear-repair-obstruction-on-the-actual-closed-bundle}{%
\subsection{4. A nonlinear repair obstruction on the actual closed
bundle}\label{a-nonlinear-repair-obstruction-on-the-actual-closed-bundle}}

\hypertarget{the-closed-detector-and-the-whole-fibre}{%
\subsubsection{4.1. The closed detector and the whole
fibre}\label{the-closed-detector-and-the-whole-fibre}}

Let \(p\) be the fibre genus. The global TT tensor \(Q_\varepsilon\),
from a concentrated base quadratic differential, satisfies \[
\beta^{-1}\!\int\langle E(G_\varepsilon),Q_\varepsilon\rangle
=-c_0\varepsilon^3(1+o(1)),\quad c_0=\frac1{2\pi^2},\quad E(g)=\operatorname{Ric}_g+g,
\] \[
\beta^{-1}\|J_GQ_\varepsilon\|_1\le C\varepsilon^2,\qquad
\beta^{-1}\!\int|Q_\varepsilon|_*\longrightarrow M_0=8\pi^2(p-1).
\tag{4.1}
\] The nuclear norm is dual to the operator norm. The baseline Taylor
estimate bounds arbitrary quadratic responses but does not give their
sign. Base-TT quadratic responses vanish; other perturbations can cancel
this single pairing while enlarging the full defect.
\href{../proofs/construction.md}{Round 1} retains that distinction.

Round 5 passes to the whole finite-area fibre limit, preserving both
cusps and the thick part. A formal double-cylinder limit could lose
detector mass. Classifying the bounded whole-fibre Jacobi response
proves its quadratic pairing nonpositive, hence for every
\(h_\varepsilon=O_{C^2}(\varepsilon^{3/2})\), \[
\liminf\varepsilon^{-3}\delta_-(G_\varepsilon+h_\varepsilon)
\ge\frac1{16\pi^4(p-1)}.
\] No symmetry, support, or gauge condition is imposed. See
\href{../proofs/round5_construction.md}{the whole-fibre proof} and
\href{../proofs/round5_receiver_independent.md}{independent receiver
derivation}.

\hypertarget{a-uniform-inverse-in-the-receivers-measure}{%
\subsubsection{4.2. A uniform inverse in the receiver's
measure}\label{a-uniform-inverse-in-the-receivers-measure}}

On the comparison product of the base cylinder and the actual closed
pinching fibre \(F_\varepsilon\), after angular averaging, set \[
d\nu=\operatorname{sech}t\,dt\,dA_{F_\varepsilon},\qquad
K=f(aT+bS)+A+J(t)B,\quad A,B\in\mathrm{TT}(F_\varepsilon).
\] The uniform weighted estimate is \[
V=K+r,\qquad \|r\|_{H^2_\nu}\le C\|LV\|_{L^2_\nu}.
\tag{4.2}
\] It uses the scalar gap of the specified nonseparating one-node
degeneration: the thick complement remains connected. A separating pinch
need not obey that gap. Off fibre TT tensors the fibre operator has gap
two. Horizontal resonances are handled by \[
(-\partial_t+\tanh t)(\partial_t+2\tanh t)=-P_t+2-4f,
\] whose reversed product is the coercive mixed operator \(-P_t+2-f\).
The vertical TT equation retains both solutions \(A+JB\), instead of
discarding transfers as nonintegrable in a different measure.

For every kernel element, \[
\mathcal B(K,K):=\frac12\int
\langle D^2E[K,K],Q_0\rangle
=-\frac{\pi}{8}\|B\|_{L^2(F_\varepsilon)}^2\le0.
\tag{4.3}
\] A determinant-preserving fibre path gives the sign: its horizontal
Ricci variation is minus the squared fibre velocity. The bilinear
estimate \(|\mathcal B(U,V)|\le C\|U\|_{H^2_\nu}\|V\|_{H^2_\nu}\)
controls the positive response by \(C\|LV\|(\|V\|+\|LV\|)\). This is not
a uniform maximum-norm inverse or a parametrization of all global bundle
deformations. See \href{../proofs/round6_weighted_inverse.md}{the
weighted proof}.

\hypertarget{arbitrary-closed-perturbations-and-the-linear-size-ball}{%
\subsubsection{4.3. Arbitrary closed perturbations and the linear-size
ball}\label{arbitrary-closed-perturbations-and-the-linear-size-ball}}

Put \(a=\|h\|_{C^2(G_\varepsilon)}\) and
\(\eta=\delta_-(G_\varepsilon+h)+\varepsilon^2+\varepsilon a+a^2\). The
actual base collar is a marked product with angular period \(\beta\);
twist clutching occurs at a radial face, not at angular periodicity.
Cutoff, averaging, and product comparison produce forcing \(O(\eta)\).
The oscillatory part has \(\|V_{\rm osc}\|_{H^1_\nu}^2\le C\beta a^2\).
Boundary flux is estimated, not set to zero.

Gauge removal has a nonzero forcing term. If \(H=V+\mathcal L_Xg_*\) and
\(F=DE(H)\), naturality gives \[
D^2E[H,H]-D^2E[V,V]=2\mathcal L_XF-DE[W].
\] The first term costs \(Ca\eta\); only the second disappears against
the exact detector. Including it, background comparison, exterior
detector mass, and the cubic remainder gives \[
\begin{split}
\beta^{-1}\!\int\langle E(G+h),Q_\varepsilon\rangle
\le -c_0\varepsilon^3(1+o(1))
+C[\varepsilon^2a+a\eta+\eta^2+(\varepsilon+\beta)a^2+a^3].
\end{split}
\tag{4.4}
\] Choose \(a\le c_*\varepsilon\) with sufficiently small fixed
\(c_*>0\). The error is then smaller than half the source, and (4.1)
gives \[
\boxed{\ \|h\|_{C^2(G_\varepsilon)}\le c_*\varepsilon
\Longrightarrow\delta_-(G_\varepsilon+h)\ge c_1\varepsilon^3.\ }
\tag{4.5}
\] For \(h=o(\varepsilon)\), the coefficient remains \(c_0/M_0\). After
reducing the radius by a fixed factor, rescaling excludes all Einstein
constants in this ball. The inherited \(O(\varepsilon^2)\) correction
achieves cubic order: (4.5) is an optimal-order statement in this
neighborhood. See \href{../proofs/round6_construction.md}{the closed
proof} and \href{../proofs/round6_receiver_review.md}{review}.

A uniform \(C^0\) modulation inverse is false. A global TT field
\(U_\varepsilon=a_\varepsilon Q_\varepsilon\), varying across the fibre
collar, has unit maximum size, residual \(\Theta(1/|\log\varepsilon|)\),
and stays a fixed distance from reference tangents plus arbitrary
infinitesimal diffeomorphisms. This
\href{../proofs/round5_modulation.md}{counterexample} explains the
receiver measure. Its much larger residual does not defeat (4.5).
Neither argument decides large \(C\varepsilon\), arbitrary \(o(1)\), or
distant repairs. At \(a=C\varepsilon\), several cubic-scale terms in
(4.4) compete, and their combined sign is not established.

\hypertarget{charged-local-geometry-and-complete-one-ended-charge-release}{%
\subsection{5. Charged local geometry and complete one-ended charge
release}\label{charged-local-geometry-and-complete-one-ended-charge-release}}

The baseline already proves stable Koiso Einstein germs with nonzero
conserved charge, all Cauchy constraints, and periodic torus data. Round
6 supplies an explicit parametrization, not first charged existence. It
uses \(M(\theta)=R_\theta\operatorname{diag}(e^a,e^{-a})R_\theta^T\),
\(Q_0=q^TM^{-1}q/2\), and \(b=2\sinh a/\sqrt{12+2Q_0}\). The analytic
\(\operatorname{Ric}=-3g\) evolution has central \(H=0\), \(H'=3-Q_0\),
positive orbit-plane curvature \(3+Q_0/2\), and a strictly stable center
when \(\max Q_0<3\). It is a local analytic germ, not smooth Cauchy
well-posedness or continuation to complete ends. Compare
\href{../baseline/research/20260916_global.tex}{the baseline
charged-center subsection} with \href{../proofs/round6_global.md}{the
explicit refinement}.

A separate complete construction releases charge with small full Ricci
error. For fixed smooth \(\chi=1\) on \([0,1]\), \(\chi=0\) on
\([2,\infty)\), put \[
m_R(r)=\chi(r/R),\quad V_R=r^2-m_R/r,\quad
g_R=\frac{dr^2}{V_R}+V_Rd\vartheta^2+r^2(dx^2+dy^2),\quad r\ge1.
\] Period \(4\pi/3\) for \(\vartheta\) makes \(r=1\) a smooth bolt. The
topology is \(D^2\times T^2\); \(r\le R\) is exact toral Einstein
black-hole geometry and \(r\ge2R\) is exactly hyperbolic. The
Ricci-error eigenvalues are \[
m_R''/(2r),\quad m_R''/(2r),\quad m_R'/r^2,\quad m_R'/r^2,
\] so the global defect is \(O(R^{-3})\). For integral torus generators
\(\partial_x+c\partial_\vartheta,\partial_y\), the charge is
\((c(3m_R-rm_R'),0)\); small mixed Ricci error changes it through the
transition. This \href{../proofs/round6_global.md}{one-ended
construction} and its concentration estimates do not provide a two-ended
torus connector.

\hypertarget{complete-two-ended-approximation-with-fixed-circle-topology}{%
\subsection{6. Complete two-ended approximation with fixed circle
topology}\label{complete-two-ended-approximation-with-fixed-circle-topology}}

\hypertarget{explicit-doubled-extremal-kottler-geometry}{%
\subsubsection{6.1. Explicit doubled extremal Kottler
geometry}\label{explicit-doubled-extremal-kottler-geometry}}

Fix a closed curvature-\(-1\) surface \((\Sigma,h_\Sigma)\) and a
\(2\pi\)-periodic circle. On \(\mathbb R_u\times S^1_\phi\times\Sigma\),
put \[
y=d\cosh u,\qquad
g_d=\frac{1+y}{3+y}\,du^2+
\frac{y^2(y+3)}{3(y+1)}\,d\phi^2+\frac{(1+y)^2}{3}h_\Sigma .
\tag{6.1}
\] For \(0<d<1/2\), it is smooth and complete, with free circle action
and two smooth AH ends of fixed conformal metric \(d\phi^2+h_\Sigma\).
Its exact global defect is \[
\|\operatorname{Ric}_{g_d}+3g_d\|_{\infty,\mathrm{op}}
=\frac{d(2d+3)}{(1+d)^2}<3d.
\tag{6.2}
\] Indeed, set \(r=(1+y)/\sqrt3\), \(F=r^2-1+2/(3\sqrt3r)\),
\(q=F(1-d^2/y^2)\). Then \(g_d=dr^2/q+Fd\phi^2+r^2h_\Sigma\), doubled at
\(y=d\). The radial and tangential Ricci errors are exactly \[
-\frac{3d^2}{y(1+y)^2},\qquad
\frac{d^2(2y+3)}{y(1+y)^2}\quad\text{(multiplicity three)}.
\] The latter dominates and decreases with \(y\), proving the continuous
supremum, not a sampled maximum. Replacing \(d^2/y^2\) by
\(d^2\chi(y)/y^2\), supported below \(y=2\) and unchanged below \(y=1\),
makes both far ends exactly extremal Kottler. The fixed transition error
is \(O(d^2)\), the center error \(O(d)\), and the physical error-support
length \(O(\log(1/d))\). Curvature and its fixed-order derivatives stay
bounded; the circle radius \(\sim d\) collapses. See
\href{../proofs/round10_kottler.md}{the full calculation}.

This realizes optical ground-state localization geometrically. For
\(w=b/a=(1+y)^{3/2}/(y\sqrt{y+3})\), the radial Neumann index operator
has positive zero mode \(w\). Normalized by its maximum, the sum of its
boundary values is \(2d\sqrt{d+3}/(1+d)^{3/2}\to0\), while error times
support length tends to zero. Thus an error-length gap without
ground-state overlap is false. The
\href{../proofs/round9_connector.md}{exact connector theorem} excludes
exact two-ended fillings in this free-circle class. It does not exclude
nonsymmetric fillings, nor is this an OT \(T^3\) connector: the boundary
is \(S^1\times\Sigma\). Boundary symmetry cannot fill the missing
hypothesis of an isometry- extension theorem; its relative
fundamental-group condition fails here because the conformal boundary is
disconnected.

\hypertarget{doubling-a-surface-cusp-on-any-fixed-circle-bundle}{%
\subsubsection{6.2. Doubling a surface cusp on any fixed circle
bundle}\label{doubling-a-surface-cusp-on-any-fixed-circle-bundle}}

Let \(P_m\to\Sigma\) have Chern number \(m\in\mathbb Z\). Suppose an
invariant exact \(\operatorname{Ric}=-3g\) source has one PE end and one
cusp with model \[
G_0=\tfrac13dr^2+c_0^2e^{-2r}\theta^2+h,\quad
h=\tfrac13h_\Sigma,\quad d\theta=f\,dA_h,\quad
f\operatorname{Area}(h)=2\pi m,
\] and relative covariant derivatives through order two of \(G-G_0\) are
\(O(e^{-\beta r})\). Li--Liu's theorem supplies such sources for every
\(m\); that source theorem constructs one PE end and one cusp. Join two
copies through \[
g_C=dt^2+a(t)^2\theta^2+h,\qquad a=d\cosh(\sqrt3t).
\] Its error eigenvalues are \(0,a^2f^2/2,-a^2f^2/2,-a^2f^2/2\); all
mixed components vanish. Match where \(a(T)=a_0\). On a fixed-width
overlap the relative two-jet difference from the outward exponential
cusp is \(O(d^2/a_0^2)\), while the source error is \(O(a_0^\beta)\). A
cutoff and the exact connection-difference formula give \[
\delta_{-3}(g_{a_0,d})\le C(a_0^\beta+a_0^2+d^2/a_0^2).
\tag{6.3}
\] Taking \(d=a_0^2\) produces complete approximate metrics on the same
\(\mathbb R\times P_m\), with both PE ends unchanged and conformal data
fixed. The connection is global on \(P_m\); no global potential on
\(\Sigma\) is assumed. Circle collapse suppresses \(a^2|F|^2\) while the
negatively curved quotient surface remains fixed. The
\href{../proofs/round10_cusp_doubling.md}{gluing proof} uses nonlinear
relative estimates, not an inverse depending on injectivity radius. The
same substitution does not handle flat \(T^3\) AH or ACH Nil cusps:
additional collapsing directions change their curvature balance.

\hypertarget{fixed-closed-metrics-from-sheared-product-blocks}{%
\subsection{7. Fixed closed metrics from sheared product
blocks}\label{fixed-closed-metrics-from-sheared-product-blocks}}

\hypertarget{the-object-and-theorem}{%
\subsubsection{7.1. The object and
theorem}\label{the-object-and-theorem}}

Take finitely many oriented compact surfaces \(B_v\) with nonempty
boundary and \(\chi(B_v)<0\), and closed surfaces \(F_v\) of common
genus \(p\ge2\). Pair every block boundary in a connected graph by \[
\Psi_e(\phi,x)=(-\phi+u_e(x),f_e(x)),
\tag{7.1}
\] where \(u_e:F_v\to\mathbb R/2\pi\mathbb Z\) and the
orientation-preserving \(f_e\) are fixed. Mapping classes need not be
finite-order, commuting, or extendible over the blocks. The resulting
closed smooth \(M\) is fixed.

\textbf{Closed gluing theorem.} There are smooth metrics \(G_\ell\) with
\[
\|\operatorname{Ric}_{G_\ell}+3G_\ell\|_{\infty,\mathrm{op}}
\le C(|\log\ell|^{-1}+\ell)\to0,\qquad
\operatorname{Vol}(G_\ell)=\frac{4\pi^2}{9}\chi(M).
\tag{7.2}
\] All fixed derivatives of the defect have the same order bound;
curvature and its fixed derivatives stay bounded. Constants depend on
the finite marked data, not \(\ell\). Scaling by three gives
\(e_-(M)=0\) and volume \(4\pi^2\chi(M)\). The
\href{../proofs/round11_closed_gluing.md}{proof} uses the following
separated operations.

\hypertarget{mapping-classes-and-integral-shear-on-disjoint-supports}{%
\subsubsection{7.2. Mapping classes and integral shear on disjoint
supports}\label{mapping-classes-and-integral-shear-on-disjoint-supports}}

Give both surface factors curvature-\(-3\) metrics and prescribe every
base boundary length to be \(\ell\). The collar lemma supplies \[
dt^2+a_\ell(t)^2d\phi^2,\qquad
a_\ell=\frac{\ell}{2\pi}\cosh(\sqrt3t),\qquad L_\ell\asymp|\log\ell|.
\] On a signed collar \([-L_\ell,L_\ell]\), the right product pulls back
to \(dt^2+a_\ell^2(d\phi+\beta_e)^2+f_e^*h_w\), with
\(\beta_e=-u_e^*d\phi\) closed and integral, not necessarily exact.

Choose a fixed smooth hyperbolic path \(h_v\to f_e^*h_w\), constant near
parameter endpoints, and traverse it slowly on \([-L_\ell,-1]\). There
is no connection shear there. Write \(S=\tfrac12h^{-1}\dot h\),
\(H=\operatorname{tr}S\), \(k=a_\ell'/a_\ell\),
\(U=a_\ell^{-1}\partial_\phi\). All blocks of
\(E=\operatorname{Ric}+3g\) are \[
E_{tt}=-\dot H-\operatorname{tr}S^2,\quad E_{UU}=-kH,\quad
E_{FF}^{\sharp}=-D_tS-(k+H)S,\quad
E_{tF}=\operatorname{div}S-dH,\quad E_{tU}=E_{UF}=0.
\tag{7.3}
\] The path is a fixed compact smooth family:
\(S,H,\operatorname{div}S,dH=O(L_\ell^{-1})\),
\(D_tS,\dot H=O(L_\ell^{-2})\), and \(|k|\le\sqrt3\). No common
pointwise area form is needed. Angular invariance prevents
inverse-circle-radius losses.

After the path finishes, use \(0\le t\le1\) for
\(g=dt^2+a_\ell^2(d\phi+z(t)\beta_e)^2+h\), keeping \(h=f_e^*h_w\)
fixed. For \(b=\beta_e\), \(D=\operatorname{div}_h b\), and a fixed
smooth step \(z\), \[
\begin{aligned}
E_{tt}&=-a^2(z')^2|b|^2/2,& E_{UU}&=a^2(z')^2|b|^2/2,\\
E_{XY}&=-a^2(z')^2b(X)b(Y)/2,& E_{tU}&=az'D/2,\\
E_{XU}&=-a(z''+3kz')b(X)/2,& E_{tX}&=0 .
\end{aligned}
\tag{7.4}
\] The \(3k\) coefficient comes from the divergence of \(a^3F\). Since
\(a=O(\ell)\) on this fixed interval, the full operator error is
\(O(\ell)\), including the mixed entries. Outside it the connection is
flat; the right endpoint is exactly the pullback by the actual
circle-valued gluing, not an invented real primitive.

\hypertarget{smoothness-volume-curvature-and-symplectic-forms}{%
\subsubsection{7.3. Smoothness, volume, curvature, and symplectic
forms}\label{smoothness-volume-curvature-and-symplectic-forms}}

Profiles are constant near every interface, so all metric jets match.
There are finitely many disjoint active collars. Their maximum estimates
prove (7.2). Shear preserves determinant; every hyperbolic fibre metric
has the same total area. Gauss--Bonnet therefore gives exactly \[
\operatorname{Vol}(G_\ell)
=\sum_v\operatorname{Area}(B_v)\operatorname{Area}(F_v)
=\frac{4\pi^2}{9}\sum_v\chi(B_v)\chi(F_v)
=\frac{4\pi^2}{9}\chi(M).
\] Pointwise constancy of fibre area is unnecessary for this identity.

Let \(\mathcal R_0\) be the algebraic tensor of two curvature-\(-3\)
planes in the actual local splitting, and \(P\) its plane projection.
For \(\epsilon_\ell=L_\ell^{-1}+\ell\), the Gauss, Codazzi, and circle
formulas give \[
|\nabla^j(\mathrm{Rm}-\mathcal R_0)|\le C_j\epsilon_\ell,\qquad
|\nabla^jP|\le C_j\epsilon_\ell\quad(j\ge1).
\] The only mixing terms are \(S=O(L_\ell^{-1})\) or
\(az'\beta_e=O(\ell)\). Consequently \[
-3-C\epsilon_\ell\le\sec\le C\epsilon_\ell,\qquad
\|\nabla^j\mathrm{Rm}\|_\infty\to0\quad(j\ge1).
\tag{7.5}
\] Neither diameter nor injectivity radius is controlled; almost
parallel curvature does not imply a locally symmetric metric on the
closed manifold.

The metrics may also be made almost Kähler in both orientations. Isotope
each \(f_e\) to preserve the chosen fibre area forms; a boundary isotopy
extends over a collar, so the smooth type stays fixed. For a path with
equal endpoint areas \(\omega\), choose \(d\sigma_z=\omega_z-\omega\) by
a fixed linear primitive operator. Then \(\sigma_0=\sigma_1=0\). Moser's
argument in an auxiliary parameter on \((1-s)\omega+s\omega_z\) produces
a common-area path with identity pullback maps at both endpoints. This
avoids a spurious endpoint holonomy. The forms \[
\Omega_\pm=a_\ell dt\wedge(d\phi+z\beta_e)\pm\omega
\] match the product forms on the blocks, are closed, and have
\(\Omega_\pm^2=\pm2a_\ell dt\wedge(d\phi+z\beta_e)\wedge\omega\ne0\).
They are compatible with the two orthogonal metric planes. Their complex
structures are not asserted integrable.

\hypertarget{which-smooth-objects-are-genuinely-beyond-the-baseline}{%
\subsubsection{7.4. Which smooth objects are genuinely beyond the
baseline}\label{which-smooth-objects-are-genuinely-beyond-the-baseline}}

If all \(f_e\) are identities, projection to a common \(F\) survives:
this is an old multitwist surface-bundle mechanism after interchanging
the two surface roles. Finite isometry holonomy also requires care.
Instead self-glue \(F_{h,2}\times\Sigma_g\), \(h\ge1,g\ge2\), with
nonzero \(u\) and infinite-order \(f\). The group has presentation \[
\langle A\times\Gamma,t\mid tc_1t^{-1}=c_2^{-1},\
t\gamma t^{-1}=c_2^{u(\gamma)}f_*(\gamma)\rangle .
\] The \href{../proofs/round11_graph_topology.md}{group proof} rules out
every extension of one closed hyperbolic surface group by another.
Asphericity and positive Euler characteristic exclude the remaining
oriented closed surface-bundle alternatives. These examples are not
merely renamed \(E_D\)'s. They are also not the infranil-fibre higher
graph manifolds of existing Einstein obstruction theorems: their walls
have nonamenable group \(\mathbb Z\times\Pi_g\), and their pieces are
products of hyperbolic surfaces.

\hypertarget{what-exact-repair-of-the-new-closed-inputs-must-change}{%
\subsection{8. What exact repair of the new closed inputs must
change}\label{what-exact-repair-of-the-new-closed-inputs-must-change}}

On a circle-invariant collar write \[
g=dt^2+h_t+a^2\theta^2,\quad S=\tfrac12h^{-1}\dot h,\quad H=\operatorname{tr}S,
\quad d\theta=dt\wedge\beta+B,\quad Q=|\beta|^2+\tfrac14|B|_{\rm tens}^2.
\] The exact vertical-minus-normal Ricci identity is \[
D=\operatorname{Ric}(U,U)-\operatorname{Ric}(T,T)
=H'-(a_t/a)H+|S|^2-\Delta_ha/a+a^2Q.
\tag{8.1}
\] For an Einstein metric \(D=0\), irrespective of its constant. When
\(a=a(t)\), its positive receiver is \[
\int_{-L}^L(|S|^2/a+aQ)\,dt=H(-L)/a(-L)-H(L)/a(L).
\tag{8.2}
\] Let \(I_a=\int a\,dt\), and let \(C_\partial\) be the maximum norm of
the right side. Logarithmic length variation and Cauchy--Schwarz give
\(d_{\mathcal T}([h_-],[h_+])^2\le I_aC_\partial\). For radii comparable
within factor \(\Lambda\) to the cosh radius, and endpoint
\(|H|\le\eta\), \[
\eta\ge\frac{\sqrt3}{4\Lambda^2}d_{\mathcal T}([h_-],[h_+])^2.
\tag{8.3}
\] The collapsing scale cancels. Pseudo-Anosov gluing has translation
length \(\log\lambda(f)>0\), even if the reference conformal structure
varies. This parameter-independent obstruction concerns an invariant,
unit-normal-foliation sector, unlike the arbitrary-tensor result of §4.

Spatial radius variation is a real escape:
\(\int_F\Delta a/a^2=2\int_F|\nabla a|^2/a^3\), not zero. In the
common-area sector \(H\equiv0\), multiplying (8.1) by \(a\) and
integrating over the closed fibre forces \(S=d\theta=0\); then
\(\Delta a=0\), and the exact metric is locally a product of
equal-curvature surfaces. See \href{../proofs/round12_repair.md}{the
full repair proof}.

The separate \href{../proofs/round12_symmetry.md}{global symmetry
theorem} says that, for these self-glued examples with \(u\ne0\), every
metric has discrete universal-cover isometry group and every Einstein
metric has no nonzero local Killing germ anywhere. Any exact solution
must therefore break circle symmetry on every open neck, not merely
abandon \(a=a(t)\). This qualitative theorem is not a distance bound
from asymmetric solutions.

\hypertarget{applications-unsuccessful-connections-and-the-global-gap}{%
\subsection{9. Applications, unsuccessful connections, and the global
gap}\label{applications-unsuccessful-connections-and-the-global-gap}}

The inherited free-quotient mechanism realizes independent twists and
the noncommuting full-generator bundle \(E_{\rm full}\). The
\href{../proofs/round8_multitwist.md}{multitwist study} adds exact
representation calculations on those same objects: every rank-one
representation of \(E_{\rm full}\) factors through the base. This is not
a new slow-interpolation mechanism or an arbitrary Einstein obstruction.

The log-Weyl route also acquired a selection rule. Under pointwise
approximation, exterior convergence, and vanishing replacement volume,
\href{../proofs/round7_endpoint.md}{the concentration theorem} gives \[
T=T_0+E/(2\pi^2),\qquad
\alpha^2(M)\le\min\{V_0/(2\pi^2),J_0+E/(12\pi^2)\},
\] where \(E\ge0\) is concentrated self-dual Weyl energy and
\(T=2\chi+3\tau\). Thus \(T-\tfrac23(J_0+E/(12\pi^2))=F_0+4E/(9\pi^2)\),
\(F_0\ge0\). Concentration cannot simply move locally symmetric bulk
into the critical monopole range. The same special-monopole hypotheses
cannot simultaneously violate the separately established near-Einstein
gap. This is a restriction on compatible inputs and obstructions, not a
new metric.

Two failed shortcuts remain relevant. Complete-model face flux cannot be
discarded from compact convergence: shrinking physical periods and the
exterior receiver matter. Growing AH cores also cannot be inserted into
fixed closed topology without scale and volume matching. The
construction in §7 solves a marked compact gluing problem instead of
invoking that unproved AH-to-cusp closure. Classical Kodaira fibrations
supply \href{../proofs/round10_bundle.md}{exact Einstein control
families} with infinite monodromy and normalized harmonic gap tending to
zero. Their positive signature differs from \(E_D\); they refute an
overbroad strategy, not the zero-signature problem.

The final closed inputs include manifolds outside all oriented surface
bundles, with \(e_-(M)=0\), bounded curvature, and almost parallel
curvature. Separate topology excludes nonnegative Einstein constants and
several special negative geometries. The remaining question is the
existence or nonexistence of a completely nonsymmetric negative Einstein
metric on those same smooth manifolds. Closing it would settle a
negative-branch instance, not both signs of the project's disjunctive
Yau26 formulation.

\hypertarget{computation-and-proof-boundaries}{%
\subsection{10. Computation and proof
boundaries}\label{computation-and-proof-boundaries}}

The source notes include independent coordinate Ricci calculations for
the cylinders, all tensor blocks, Kottler family, circle-cusp core, and
both regions of the closed gluing. The elliptic inverses, analytic
continuation, topology, and existence arguments are mathematical proofs,
not conclusions certified by a finite script. Checks reject missing
trace terms, incorrect mixed coefficients, and omitted gauge responses.
Their scope and historical execution appear in
\href{../proofs/round6_verification.md}{Round 6},
\href{../proofs/round10_verification.md}{Round 10},
\href{../proofs/round11_verification.md}{Round 11}, and
\href{../proofs/round12_verification.md}{Round 12}. This chapter
reorganizes those results; it does not report a fresh execution of every
program. Its coverage table identifies each supporting proof, review,
refinement, and superseded intermediate route.

\end{document}
